\documentclass[11pt,a4paper]{article}
\usepackage{ftm}

\title{Programas de Pesquisa FTM\\
       para os Seis Problemas do Milênio Abertos\\[6pt]
       \large Núcleos Rígidos, Lacunas Formais e Obrigações de Prova}
\author{Federated Trivalent Manifold Research Programme}
\date{2026}

\begin{document}
\maketitle

\begin{abstract}
Este documento reúne os seis cadernos de lacunas formais para os
Problemas do Milênio ainda abertos (Hipótese de Riemann, P vs NP,
Birch–Swinnerton-Dyer, Hodge, Navier–Stokes, Yang–Mills e gap de massa),
elaborados no quadro do programa \emph{Federated Trivalent Manifold} (FTM).

Para cada problema são apresentados:
(i)~o enunciado oficial Clay;
(ii)~os principais programas de pesquisa em curso e onde o formalismo
FTM se encaixa;
(iii)~as obrigações formais numeradas que devem ser satisfeitas
antes de qualquer alegação forte;
(iv)~as lacunas identificadas com status \OPEN.

\textbf{Declaração de estado científico.}
Nenhuma seção constitui solução de nenhum problema.
Uma conexão só é tratada como matemática quando há
definições, enunciado formal e prova auditável;
o restante é rotulado como programa, conjectura ou analogia.
As formulações oficiais são as do
\href{https://www.claymath.org/millennium-problems/}{Clay Mathematics Institute}.
\end{abstract}

\tableofcontents
\clearpage

%% ---------------------------------------------------------------
\section{Estrutura comum e metodologia}
%% ---------------------------------------------------------------

\subsection{Rigidez local versus defeito global}

Todos os seis problemas compartilham a mesma arquitetura:

\begin{center}
\begin{tabular}{lll}
\toprule
Problema & Rigidez local (fácil) & Defeito global (difícil) \\
\midrule
Riemann   & $\Delta_{\mathrm{par}}=\delta^2$ geométrico
          & $\Drad=\sum w$ convergência uniforme \\
P vs NP   & Verificação polinomial local
          & Busca polinomial global \\
BSD       & Fatores $L_p(E,s)$ locais
          & Posto de Mordell–Weil global \\
Hodge     & Descida e tipo Hodge locais
          & Sobrejetividade do mapa de ciclos \\
Navier–Stokes & Regularidade local $L^3$
              & Estimativa crítica uniforme global \\
Yang–Mills & Regularização discreta $a>0$
           & Gap uniforme no limite $a\to0$ \\
\bottomrule
\end{tabular}
\end{center}

O manifesto FTM — \emph{``matemática finita não alcança o infinito
sem ponte de conservação''} — é estruturalmente correto para todos.
Mas a ponte requer definições precisas:
operador $T$ com $Q_W=\|Tg\|^2$, erros $e_N\to0$, ruído $\omega\to0$;
não basta analogia qualitativa.

\subsection{Metodologia Aceitar-Núcleo}

Para cada problema:
\begin{enumerate}[leftmargin=2em]
\item \textbf{Isolar.} Transcrever enunciado Clay, objetos,
  quantificadores, alternativas e o que não basta.
\item \textbf{Núcleo.} Identificar obstrução mínima (tabela acima).
\item \textbf{Aceitar hipótese de trabalho.}
  Testar circularidade: a hipótese $H$ continua válida se a
  conjectura for falsa? Se não, marcar EQUIVALENTE.
\item \textbf{Derivar.} Sob $H$, construir algoritmos ou funcionais.
\item \textbf{Testar adversarialmente.}
  Contramodelo sintético (zeros fora da reta, circuitos que quebram
  verificação, dados NS quase blow-up).
\item \textbf{Classificar.}
  \PROVED\ / \CONDITIONAL\ / \OPEN\ para cada afirmação.
\end{enumerate}

\subsection{Escala de maturidade}

\begin{description}[leftmargin=4em,style=nextline]
\item[M0] Nenhum objeto FTM relevante construído;
  as obrigações são apenas enunciadas.
\item[M1] Objetos construídos conceitualmente;
  obrigações enunciadas e parcialmente testadas em exemplos.
\item[M2] Pelo menos uma obrigação provada;
  exemplos não triviais verificados.
\item[M3] Cadeia de implicações com ao menos um elo publicado.
\end{description}

\clearpage

%% ---------------------------------------------------------------
\section{Hipótese de Riemann}
%% ---------------------------------------------------------------

\textit{Caderno completo:} \texttt{riemann.tex}.

\medskip
\noindent\textbf{Enunciado.}
$\zeta(s)=0$, $0<\Re s<1$ $\Rightarrow$ $\Re s=\tfrac12$.

\medskip
\noindent\textbf{Lacunas principais.}
\begin{gap}
  Operador $D_F$ auto-adjunto com $\det_{\mathrm{reg}}(D_F-z)=E(z)\xi(\tfrac12+iz)$
  (obrigações RH1–RH4) não foi construído. \OPEN
\end{gap}
\begin{gap}
  $\sum_\rho|\rho|^{-2}$ não-holomorfo: convergência uniforme de $\Drad(T)$
  sem assumir RH não foi provada. \OPEN
\end{gap}
\begin{gap}
  $G_N\succeq0$ com $e_N\to0$ e $\omega\to0$ simultâneos não foi estabelecido. \OPEN
\end{gap}

\MATURITY{M0. Programa off-critical (Eixo~E) é o entregável mais próximo.}

\clearpage

%% ---------------------------------------------------------------
\section{P versus NP}
%% ---------------------------------------------------------------

\textit{Caderno completo:} \texttt{p-vs-np.tex}.

\medskip
\noindent\textbf{Enunciado.}
$P\stackrel{?}{=}NP$. Verificação polinomial implica busca polinomial?

\medskip
\noindent\textbf{Lacunas principais.}
\begin{gap}
  Barreiras de relativização, natural proofs e algebrização não foram
  transpostas por nenhuma técnica FTM. \OPEN
\end{gap}
\begin{gap}
  A linguagem FTM-DESCENT carece de codificação finita definitiva (C1 aberta);
  sem C1, C2–C4 não são formuláveis. \OPEN
\end{gap}
\begin{gap}
  Mesmo NP-completude de FTM-DESCENT não separaria $P$ de $NP$. \OPEN
\end{gap}

\MATURITY{M0. C1–C4 todas abertas.}

\clearpage

%% ---------------------------------------------------------------
\section{Birch–Swinnerton-Dyer}
%% ---------------------------------------------------------------

\textit{Caderno completo:} \texttt{bsd.tex}.

\medskip
\noindent\textbf{Enunciado.}
$\ord_{s=1}L(E,s)=\rank E(\mathbb Q)$ com fórmula refinada
para o coeficiente líder.

\medskip
\noindent\textbf{Lacunas principais.}
\begin{gap}
  Caso $r\ge2$: não há método provado para construir pontos racionais
  independentes via teoria analítica de $L(E,s)$. \OPEN
\end{gap}
\begin{gap}
  Main Conjecture de Iwasawa em generalidade plena: igualdade de ideais
  característicos não provada para todos os primos e reduções. \OPEN
\end{gap}
\begin{gap}
  Funtor $\Phi:\mathrm{Ell}_{\mathbb Q}\to\mathrm{FTM}_{\mathrm{arith}}$
  satisfazendo BSD1–BSD5 não foi construído. \OPEN
\end{gap}

\MATURITY{M0/M1 conceitual. BSD1–BSD5 abertas.}

\clearpage

%% ---------------------------------------------------------------
\section{Conjectura de Hodge}
%% ---------------------------------------------------------------

\textit{Caderno completo:} \texttt{hodge.tex}.

\medskip
\noindent\textbf{Enunciado.}
Para $X/\mathbb C$ projetiva não singular, toda classe de Hodge é
combinação $\mathbb Q$-linear de classes de ciclos algébricos.

\medskip
\noindent\textbf{Lacunas principais.}
\begin{gap}
  Conjecturas padrão de Grothendieck ($A$–$D$) abertas; implicariam Hodge. \OPEN
\end{gap}
\begin{gap}
  H-SURJ (sobrejetividade do mapa de ciclos) é a lacuna central;
  H-DESC e H-TYPE provados não a implicam. \OPEN
\end{gap}
\begin{gap}
  Extensão de Cattani–Deligne–Kaplan para variedades gerais de dimensão $\ge3$
  não existe. \OPEN
\end{gap}

\MATURITY{M0. H-DESC, H-TYPE, H-CYCLE, H-SURJ todas abertas.}

\clearpage

%% ---------------------------------------------------------------
\section{Navier–Stokes}
%% ---------------------------------------------------------------

\textit{Caderno completo:} \texttt{navier-stokes.tex}.

\medskip
\noindent\textbf{Enunciado.}
Existe solução suave globalmente definida para $u_0\in\mathcal S(\mathbb R^3)$?
Ou pode ocorrer blow-up em tempo finito?

\medskip
\noindent\textbf{Lacunas principais.}
\begin{gap}
  NS4 — estimativa crítica uniforme em norma $L^3$ ou $BMO^{-1}$
  suficiente para regularidade global não foi obtida. \OPEN
\end{gap}
\begin{gap}
  Cenário Hou–Luo: nenhuma verificação rigorosa (intervalar) de blow-up
  em tempo finito para dado suave. \OPEN
\end{gap}
\begin{gap}
  Não-unicidade de soluções fracas (Buckmaster–Vicol) não implica
  breakdown da solução forte; o gap entre Leray e Prodi–Serrin
  crítico permanece aberto. \OPEN
\end{gap}

\MATURITY{M1 conceitual. NS1–NS4 abertas; NS4 é a lacuna decisiva.}

\clearpage

%% ---------------------------------------------------------------
\section{Yang–Mills e Gap de Massa}
%% ---------------------------------------------------------------

\textit{Caderno completo:} \texttt{yang-mills.tex}.

\medskip
\noindent\textbf{Enunciado.}
$SU(N)$ Yang–Mills quântico em $\mathbb R^4$ existe
e tem gap de massa $\Delta>0$.

\medskip
\noindent\textbf{Lacunas principais.}
\begin{gap}
  Extensão do programa de Balaban para 4D: controle uniforme do limite
  contínuo a todas as ordens de renormalização não provado. \OPEN
\end{gap}
\begin{gap}
  Lei de área para Wilson loop em $SU(N)$ 4D contínuo: não provada. \OPEN
\end{gap}
\begin{gap}
  YM5 — gap de massa $\Delta>0$ sobrevivendo simultaneamente a $a\to0$
  e $\Lambda\to\infty$ não provado. \OPEN
\end{gap}

\MATURITY{M1 conceitual. YM1–YM5 todas abertas.}

\clearpage

%% ---------------------------------------------------------------
\section{Sumário de lacunas e prioridades}
%% ---------------------------------------------------------------

\subsection{Tabela de status}

\begin{center}
\begin{tabular}{llll}
\toprule
Problema & Obrigação central & Status & Entregável próximo \\
\midrule
Riemann    & RH1–RH4 (operador + identidade)
           & M0, \OPEN & Programa off-critical (MoC) \\
P vs NP    & C1–C4 (codificação + barreiras)
           & M0, \OPEN & — \\
BSD        & BSD1–BSD5 (funtor + equivalência)
           & M0/M1, \OPEN & — \\
Hodge      & H-SURJ (sobrejetividade ciclos)
           & M0, \OPEN & — \\
Navier–Stokes & NS4 (estimativa crítica)
              & M1, \OPEN & Verificação intervalar \\
Yang–Mills & YM1–YM5 (construção + gap)
           & M1, \OPEN & — \\
\bottomrule
\end{tabular}
\end{center}

\subsection{Prioridade recomendada}

\begin{enumerate}[leftmargin=2em]
\item \textbf{Riemann Eixo~E (off-critical):}
  Implementar $\Drad(T)$ com aritmética de precisão arbitrária (Arb),
  publicar algoritmo de localização espectral como artigo Math.\ Comp.
  Horizonte: 6--12 meses. Probabilidade de publicação: $\sim70\%$.
\item \textbf{Navier–Stokes NS4 numérico:}
  Desenvolver verificador intervalar para cenário Hou–Luo.
  Horizonte: 1--2 anos.
\item \textbf{Yang–Mills e BSD:}
  Exigem avanços estruturais em QFT construtiva e teoria de Iwasawa,
  respectivamente. Horizonte mínimo: 5--10 anos.
\item \textbf{P vs NP e Hodge:}
  Dependem de rupturas em barreiras de complexidade (GCT)
  e na teoria de ciclos (motivos). Horizonte indeterminado.
\end{enumerate}

\subsection{O que não constitui progresso}

\begin{itemize}[leftmargin=2em]
\item Ajuste numérico de zeros ou estatísticas GUE sem RH1–RH4.
\item Linguagem FTM-DESCENT sem codificação finita (C1 aberta).
\item Gap espectral em grafo finito sem limites YM2 uniformes.
\item Classes de Hodge em exemplos sem H-SURJ geral.
\item Convergência para solução fraca sem NS4.
\item Analogia estrutural (``defeito local $\leftrightarrow$ defeito global'')
  sem implicação lógica provada.
\end{itemize}

\bigskip\noindent
\textbf{Conclusão auditável.}
Os seis programas são independentes entre si (analogia $\neq$ redução).
Nenhuma lacuna foi fechada neste estágio.
O valor deste documento é transformar seis problemas abertos em
programas falsificáveis, com entregáveis auditáveis e critérios
explícitos de interrupção, sem declarar avanço substantivo não provado.

\end{document}
